% =============================================================================
% APPENDIX INDEX: The EBF Appendix System
% =============================================================================
% LANGUAGE: english
% =============================================================================
%
% HEADER BLOCK
% -----------------------------------------------------------------------------
% Appendix:      INDEX
% Category:      REF (Reference)
% Version:       2.0 (January 2026)
% Dependencies:  All appendices, All chapters
% Primary Ch.:   All
% Secondary Ch.: ---
% -----------------------------------------------------------------------------
%
% This document provides the complete overview and navigation guide for all
% EBF appendices. It establishes the naming convention, categorization,
% and interdependencies of the 139 appendices across 24 chapters.
%
% =============================================================================

\chapter{The EBF Appendix System}
\label{app:index}

\begin{abstract}
The EBF framework is supported by 166 appendices organized into 8 categories.
This index provides the naming convention, brief descriptions, and cross-reference
structure. Each appendix follows the format: \textbf{CATEGORY-NAME: Descriptive Title}.
\end{abstract}


\begin{tcolorbox}[
    colback=orange!5!white,
    colframe=orange!75!black,
    title={\textbf{Appendix Scope: Ziel / In-Scope / Out-of-Scope / Constraints}},
    fonttitle=\bfseries
]

\textbf{Ziel (Objective):}
\begin{itemize}[nosep]
    \item Provide systematic analysis for REF integration
\end{itemize}

\textbf{In-Scope:}
\begin{itemize}[nosep]
    \item Core analysis and findings
    \item Framework integration points
\end{itemize}

\textbf{Out-of-Scope:}
\begin{itemize}[nosep]
    \item Detailed empirical replication $\rightarrow$ METHOD appendices
\end{itemize}

\textbf{Constraints:}
\begin{itemize}[nosep]
    \item Analysis bounded by available data
\end{itemize}

\end{tcolorbox}

\begin{tcolorbox}[colback=gray!5!white, colframe=gray!75!black,
    title=Appendix: INDEX | Category: REF]
\small
\begin{tabular}{@{}ll@{}}
\textbf{Version:} & 2.0 (January 2026) \\
\textbf{Dependencies:} & All appendices (166), All chapters (25) \\
\textbf{Primary Chapter:} & All chapters \\
\textbf{Purpose:} & Navigation and cross-reference guide \\
\end{tabular}
\end{tcolorbox}

%%%%%%%%%%%%%%%%%%%%%%%%%%%%%%%%%%%%%%%%%%%%%%%%%%%%%%%%%%%%%%%%%%%%%%%%%%%%%%%
\section{Naming Convention and Structure}
\label{sec:index-convention}
%%%%%%%%%%%%%%%%%%%%%%%%%%%%%%%%%%%%%%%%%%%%%%%%%%%%%%%%%%%%%%%%%%%%%%%%%%%%%%%

\begin{tcolorbox}[colback=blue!5!white, colframe=blue!75!black, 
    title=Appendix Naming Convention]
    
Every appendix follows the naming pattern:

\begin{center}
\textbf{CATEGORY-NAME: Descriptive Title}
\end{center}

\textbf{Components:}
\begin{itemize}[nosep]
\item \textbf{CATEGORY}: One of 8 prefixes indicating the appendix type
\item \textbf{NAME}: Short, memorable identifier (often a question word or acronym)
\item \textbf{Descriptive Title}: Full explanation of contents
\end{itemize}

\textbf{Example:} \texttt{CORE-WHO: The Welfare Hierarchy}
\begin{itemize}[nosep]
\item CORE = Fundamental theory category
\item WHO = Answers "Who has utility?"
\item The Welfare Hierarchy = Describes the level structure
\end{itemize}

\end{tcolorbox}

%%%%%%%%%%%%%%%%%%%%%%%%%%%%%%%%%%%%%%%%%%%%%%%%%%%%%%%%%%%%%%%%%%%%%%%%%%%%%%%
\section{The Eight Categories}
\label{sec:index-categories}
%%%%%%%%%%%%%%%%%%%%%%%%%%%%%%%%%%%%%%%%%%%%%%%%%%%%%%%%%%%%%%%%%%%%%%%%%%%%%%%

\begin{table}[htbp]
\centering
\caption{Appendix Categories Overview}
\label{tab:index-categories}
\renewcommand{\arraystretch}{1.4}
\begin{tabular}{@{}clcl@{}}
\toprule
\textbf{Prefix} & \textbf{Category} & \textbf{Count} & \textbf{Purpose} \\
\midrule
\textbf{CORE-} & Core Theory & 11 & Fundamental building blocks of EBF (10C + Hierarchy + EIT + Principles) \\
\textbf{FORMAL-} & Formalization & 19 & Mathematical foundations and proofs \\
\textbf{DOMAIN-} & Application Domains & 22 & Economics subfield applications \\
\textbf{CONTEXT-} & Context Dimensions & 6 & The $\Psi$ framework components \\
\textbf{METHOD-} & Methodology & 27 & Estimation and validation methods \\
\textbf{PREDICT-} & Predictions & 8 & Falsifiable predictions and cases \\
\textbf{LIT-} & Literature & 61 & Research integration by author \\
\textbf{REF-} & Reference & 9 & Glossaries, examples, meta-theory, skills, defense, lead management \\
\midrule
& \textbf{Total} & \textbf{159} & \\
\bottomrule
\end{tabular}
\end{table}

%%%%%%%%%%%%%%%%%%%%%%%%%%%%%%%%%%%%%%%%%%%%%%%%%%%%%%%%%%%%%%%%%%%%%%%%%%%%%%%
%%%%%%%%%%%%%%%%%%%%%%%%%%%%%%%%%%%%%%%%%%%%%%%%%%%%%%%%%%%%%%%%%%%%%%%%%%%%%%%
\section{CORE: The Eight Fundamental Questions (10C)}
\label{sec:index-core}
%%%%%%%%%%%%%%%%%%%%%%%%%%%%%%%%%%%%%%%%%%%%%%%%%%%%%%%%%%%%%%%%%%%%%%%%%%%%%%%
%%%%%%%%%%%%%%%%%%%%%%%%%%%%%%%%%%%%%%%%%%%%%%%%%%%%%%%%%%%%%%%%%%%%%%%%%%%%%%%

\begin{tcolorbox}[colback=red!5!white, colframe=red!75!black,
    title=The CORE Appendices: Foundation of EBF]

The eight CORE appendices (10C) answer the fundamental questions of welfare theory.
Together they define the complete utility space $U^L_d(\Psi)$, its calibration, the
subjective perception filter $A(\cdot)$ that transforms potential into effective utility,
the action threshold mechanism $WAX \geq \theta$ that bridges knowing and doing,
and the behavioral change journey phase $\varphi$ that locates agents in their change process.
the action threshold mechanism $WAX \geq \theta$ that bridges knowing and doing, and
the operational process that governs EBF application itself.

\begin{center}
\begin{tikzpicture}[node distance=2.5cm]
% This is conceptual - would need tikz package
\end{tikzpicture}
\end{center}

\end{tcolorbox}

\begin{table}[htbp]
\centering
\caption{CORE Appendices: The Eight Fundamental Questions (10C)}
\label{tab:index-core}
\renewcommand{\arraystretch}{1.6}
\begin{tabular}{@{}cllll@{}}
\toprule
\textbf{Code} & \textbf{Name} & \textbf{Question} & \textbf{Title} & \textbf{Defines} \\
\midrule
AAA & \textbf{CORE-WHO} & Who has utility? & The Welfare Hierarchy & Levels $L$ \\
C & \textbf{CORE-WHAT} & What is utility? & The Welfare Dimensions & Dimensions $d$ \\
B & \textbf{CORE-HOW} & How do they interact? & The Interaction Structure & Complementarity $\gamma$ \\
V & \textbf{CORE-WHEN} & When does context matter? & The Context Framework & Context $\Psi$ \\
BBB & \textbf{CORE-WHERE} & Where do the numbers come from? & Parameter Estimation & $\Theta, E(\theta)$ \\
AU & \textbf{CORE-AWARE} & How aware am I of which utility? & The Awareness Function & $A(\cdot)$ \\
AV & \textbf{CORE-READY} & How ready am I to act? & The Willingness Function & $WAX, \theta$ \\
AW & \textbf{CORE-STAGE} & Where in the journey? & The Behavioral Change Journey & Phase $\varphi$ \\
AW & \textbf{CORE-STAGE} & Where in the change journey? & The Behavioral Change Journey & $S(t), dS/dt$ \\
HI & \textbf{CORE-HIERARCHY} & How do decisions stratify? & Decision Hierarchy \& Strategic Coherence & L$_0$-L$_1$-L$_2$-L$_3$ \\
IE & \textbf{CORE-EIT} & How to intervene? & Emergent Intervention Theory & $\vec{I} \in [0,1]^9$, E-Modes \\
PR & \textbf{CORE-PRINCIPLES} & What defines EBF? & Fundamental Axioms & P1-P7, EBF $\Leftrightarrow$ $\wedge$(P1-P7) \\
\bottomrule
\end{tabular}
\end{table}

\subsection{CORE-WHO: The Welfare Hierarchy (Appendix WHO)}

\begin{tcolorbox}[colback=gray!5!white, colframe=gray!75!black]
\textbf{Question:} Who can have utility? At what levels does welfare exist?

\textbf{Answer:} Utility exists at multiple levels from Individual ($L=0$) to 
Meta (all sentient beings, $L=\infty$). Levels are not arbitrary but \textit{endogenous}---they 
exist because of shared identity (IDN) or collective utility (KNU).

\textbf{Key Content:}
\begin{itemize}[nosep]
\item The 9-level hierarchy: Individual → Family → Tribe → Organization → Network → Region → Nation → Civilization → Meta
\item Axioms SRL-1 through SRL-8 defining level structure
\item The General Theory of Welfare Hierarchies (GTWH)
\item Cross-level complementarity $\gamma^{cross}_{L,L'}$
\item Critical foundations addressing 10 disciplinary objections
\end{itemize}

\textbf{Key Formula:}
\begin{equation}
U^L = \sum_{i \in M^L} \omega^L_i \cdot U^{ACE-1}_i + \sum_{i < j} \gamma^L_{ij} \cdot U^{ACE-1}_i \cdot U^{ACE-1}_j
\end{equation}

\textbf{Lines:} 4,097 | \textbf{Sections:} 26 | \textbf{Cross-References:} CORE-WHAT, CORE-HOW, CORE-WHERE (parameters)
\end{tcolorbox}

\subsection{CORE-WHAT: The Welfare Dimensions (Appendix WAT)}

\begin{tcolorbox}[colback=gray!5!white, colframe=gray!75!black]
\textbf{Question:} What constitutes utility? What are the dimensions of welfare?

\textbf{Answer:} Utility decomposes into six FEPSDE dimensions, three categories 
(INU/KNU/IDN), four time horizons, and two valences (Gain/Pain), yielding 
144 distinct components.

\textbf{Key Content:}
\begin{itemize}[nosep]
\item The FEPSDE dimensions: Financial, Emotional, Physical, Social, Digital, Ecological
\item Meta-Axiom U-M0: Universal utility structure
\item 220 axioms: INU (59), KNU (65), IDN (95)
\item Loss aversion by dimension ($\lambda_d$)
\item The $\Gamma$-matrix: Context-welfare interaction
\item Integration with CORE-WHO (5 theoretical justifications)
\end{itemize}

\textbf{Key Formula:}
\begin{equation}
U^{i} = \sum_{d \in FEPSDE} \sum_{t=0}^{3} \delta^{i}_{d}(t) \cdot 
\left[ G^{i}_{d,t} - \lambda^{i}_{d} \cdot P^{i}_{d,t} \right]
\end{equation}

\textbf{Lines:} 4,237 | \textbf{Sections:} 15 | \textbf{Cross-References:} CORE-WHO, CORE-WHEN, CORE-WHERE (parameters)
\end{tcolorbox}

\subsection{CORE-HOW: The Interaction Structure (Appendix HOW)}

\begin{tcolorbox}[colback=gray!5!white, colframe=gray!75!black]
\textbf{Question:} How do utility components interact? What makes them complements or substitutes?

\textbf{Answer:} The complementarity parameter $\gamma$ captures non-additive 
interactions between dimensions, levels, and entities. Positive $\gamma$ indicates 
complements (synergy); negative $\gamma$ indicates substitutes (trade-offs).

\textbf{Key Content:}
\begin{itemize}[nosep]
\item Within-level complementarity $\gamma^L_{ij}$
\item Cross-level complementarity $\gamma^{cross}_{L,L'}$
\item Dimension complementarity $\gamma_{dd'}$
\item The complementarity matrix and its properties
\item Connection to Milgrom-Roberts supermodularity
\end{itemize}

\textbf{Key Formula:}
\begin{equation}
\gamma_{ij} = \frac{\partial^2 U}{\partial x_i \partial x_j} \quad
\begin{cases}
> 0 & \text{complements (synergy)} \\
= 0 & \text{independent} \\
< 0 & \text{substitutes (trade-off)}
\end{cases}
\end{equation}

\textbf{Cross-References:} CORE-WHO, CORE-WHAT, CORE-WHERE (parameters), DOMAIN-IO (Milgrom-Roberts)
\end{tcolorbox}

\subsection{CORE-WHEN: The Context Framework (Appendix CTW)}

\begin{tcolorbox}[colback=gray!5!white, colframe=gray!75!black]
\textbf{Question:} When and where does context affect welfare? How does $\Psi$ modulate behavior?

\textbf{Answer:} The 8-dimensional $\Psi$ framework captures contextual factors 
that modulate level salience, parameter values, and behavioral responses.

\textbf{Key Content:}
\begin{itemize}[nosep]
\item The 8 $\Psi$-dimensions: Economic, Social, Temporal, Spatial, Institutional, Cultural, Technological, Environmental
\item Context-dependent level salience $\alpha^L(\Psi)$
\item Parameter modulation: $\lambda(\Psi)$, $\delta(\Psi)$, $\omega(\Psi)$
\item The Awareness function $A(\Psi)$
\item Integration with EBF KON module (8904)
\end{itemize}

\textbf{Key Formula:}
\begin{equation}
L^*(\Psi) = \arg\max_L \left[ A^L(\Psi) \right]
\end{equation}

\textbf{Cross-References:} CORE-WHO (level salience), CORE-WHERE (parameters), CONTEXT-TIME, CONTEXT-SPACE
\end{tcolorbox}

\subsection{CORE-WHERE: Parameter Estimation (Appendix EST)}

\begin{tcolorbox}[colback=gray!5!white, colframe=gray!75!black]
\textbf{Question:} Where do the numbers come from? How are parameters calibrated?

\textbf{Answer:} The estimation methodology provides empirically grounded parameter
values through a hierarchy of evidence: EMP (empirical), THR (theoretical),
LLM (LLM Monte Carlo), ILL (illustrative), HYP (hypothetical).

\textbf{Key Content:}
\begin{itemize}[nosep]
\item Epistemic status tags for all parameters
\item LLM Monte Carlo estimation methodology
\item Prior specification from literature
\item Validation protocols and robustness checks
\item The fundamental trilemma: precision vs. coverage vs. verifiability
\end{itemize}

\textbf{Key Formula:}
\begin{equation}
E(\theta | \text{evidence}) = \sum_{s \in \{EMP, THR, LLM, ILL, HYP\}} w_s \cdot \hat{\theta}_s
\end{equation}

\textbf{Cross-References:} All CORE appendices (provides parameters), METHOD-LLMMC (AN)
\end{tcolorbox}

\subsection{CORE-AWARE: The Awareness Function (Appendix AWA)}

\begin{tcolorbox}[colback=gray!5!white, colframe=gray!75!black]
\textbf{Question:} How aware am I of which utility? What determines salience?

\textbf{Answer:} The Awareness function $A(\cdot) \in [0,1]$ acts as a multiplicative filter
transforming potential utility $U_{pot}$ into effective utility $U_{eff}$ at the Moment of Truth $t^*$.

\textbf{Key Content:}
\begin{itemize}[nosep]
\item The Transformation Equation: $U_{eff} = A(\cdot) \times U_{pot}$
\item 104 axioms (AWX-A1 to A90) governing awareness dynamics
\item Determinants: salience $\sigma$, cognitive capacity $K$, temporal decay $\tau$
\item Moment of Truth: Only $A(t^*)$ matters for behavior
\item Distinction from classical economics: BCM assumes $A < 1$ as norm
\end{itemize}

\textbf{Key Formula:}
\begin{equation}
U_{eff,X}(t^*) = A_X(t^*) \times U_{pot,X}(t^*)
\end{equation}

\textbf{Cross-References:} CORE-WHAT (dimensions filtered), CORE-WHEN ($\Psi$ affects $A$), CORE-READY (AWX $\to$ WAX)
\end{tcolorbox}

\subsection{CORE-READY: The Willingness Function (Appendix REA)}

\begin{tcolorbox}[colback=gray!5!white, colframe=gray!75!black]
\textbf{Question:} How ready am I to act? What bridges knowing and doing?

\textbf{Answer:} The Willingness function determines whether effective utility
$U_{eff}$ translates into action via the threshold mechanism $WAX \geq \theta$.
This is the 7th CORE question in the path from utility to behavior.

\textbf{Key Content:}
\begin{itemize}[nosep]
\item The Action Condition: $\text{Action} \Leftrightarrow WAX \geq \theta$
\item 99 axioms (WAX-A1 to A99) governing willingness dynamics
\item Thinking style parameter $\varphi \in [0,1]$ (additive vs. complementary)
\item Threshold determinants: habit, friction, social proof, regulation
\item The Intention-Behavior Gap: why knowing $\neq$ doing
\end{itemize}

\textbf{Key Formula:}
\begin{equation}
WAX = f(\varphi, U^*, \Psi) \quad \text{where } U^* = A(\cdot) \times U_{pot}
\end{equation}

\textbf{Cross-References:} CORE-AWARE (provides $U^*$), CORE-HOW ($\gamma \to \varphi$), CORE-WHEN ($\Psi$ modulates $\theta$)
\end{tcolorbox}

\subsection{CORE-STAGE: The Behavioral Change Journey (Appendix STA)}

\begin{tcolorbox}[colback=gray!5!white, colframe=gray!75!black]
\textbf{Question:} Where in the behavioral change journey is the agent located?

\textbf{Answer:} The BCJ Phase $\varphi \in \{1,2,3,4,5\}$ identifies where an individual
or collective stands in the change process---from pre-contemplation through maintenance.

\textbf{Key Content:}
\begin{itemize}[nosep]
\item The 5-Phase Model: (1) Pre-contemplation, (2) Contemplation, (3) Preparation, (4) Action, (5) Maintenance
\item Phase-dependent intervention design
\item BCJ Segments: typology of change readiness
\item Integration with WAX threshold dynamics
\item Context modulation of phase transitions
\textbf{Question:} Where in the behavioral change journey is the person? How does $S(t)$ evolve?

\textbf{Answer:} The Behavioral Change Journey (BCJ) defines the temporal dynamics
of behavioral change---how individuals move through stages from awareness to sustained action.

\textbf{Key Content:}
\begin{itemize}[nosep]
\item 85 axioms across 15 categories (BCJ-M, E, ED, D, K, S, C, B, MB, L, LF, FB, HE, TS, AP)
\item 152 integrated models from psychology, economics, and neuroscience
\item Stage position $S(t)$ and stage dynamics $dS/dt$
\item Integration with TTM, Rubicon, HAPA, COM-B, and 148 other models
\item Temporal parameters: $\tau$ (time constant), $\rho$ (persistence), $\beta$ (social contagion)
\end{itemize}

\textbf{Key Formula:}
\begin{equation}
\varphi(t) = \text{BCJ\_Phase}(t) \in \{1,...,5\}
\end{equation}

\textbf{Cross-References:} CORE-READY (WAX threshold), CORE-WHEN ($\Psi$ modulates phase), Chapter 13, 14, 15
\frac{dS}{dt} = \frac{1}{\tau} \cdot \left[ S^*(A, WAX, \Psi) - S(t) \right] + \beta \cdot \bar{S}_{peers} + \varepsilon
\end{equation}

\textbf{Cross-References:} CORE-AWARE (provides $A$), CORE-READY (provides WAX), Chapter 13 (main exposition)
\end{tcolorbox}

\subsection{CORE-EIT: Emergent Intervention Theory (Appendix IE)}

\begin{tcolorbox}[colback=gray!5!white, colframe=gray!75!black]
\textbf{Question:} What is the theoretical structure of behavioral interventions, and how does data availability constrain methodological choices?

\textbf{Answer:} Interventions exist in a continuous 9-dimensional space $[0,1]^9$ corresponding to the 10C dimensions. Each intervention emerges from the continuous space based on context analysis---not selected from a predefined catalogue. The choice of emergence mode (E-Full, E-Partial, E-None) is constrained by data availability.

\textbf{Key Content (4-Part Structure):}
\begin{itemize}[nosep]
\item \textbf{Part 1: Foundations \& Clusters (EIT-1 to EIT-3)} --- Continuous space, transformation, cluster emergence
\item \textbf{Part 2: Journey Position \& Horizon (EIT-4 to EIT-5)} --- Temporal emergence from journey phase and effect horizon
\item \textbf{Part 3: Universal Intensities (EIT-6 to EIT-8)} --- General principles for AWARE, READY, WHEN targeting
\item \textbf{Part 4: Component Functions (EIT-9 to EIT-13)} --- WHO, WHAT, HOW, WHERE, HIERARCHY intensities
\item \textbf{Master Matrix} (one-page lookup table) --- All 8 emergence functions with inputs, principles, CORE dependencies
\end{itemize}

\textbf{Additional Content:}
\begin{itemize}[nosep]
\item Three Emergence Modes: E-Full (all emerges), E-Partial (hybrid), E-None (heuristic)
\item Data Sufficiency Score $D_{\text{suff}}$ for mode selection
\item Confidence-Weighted Decision Rule (Theorem IE-T4)
\item Emergence with weak priors: structured uncertainty vs.\ false precision
\end{itemize}

\textbf{Key Formula:}
\begin{equation}
\vec{I} = (I_{\text{WHO}}, I_{\text{WHAT}}, I_{\text{HOW}}, I_{\text{WHEN}}, I_{\text{WHERE}}, I_{\text{AWARE}}, I_{\text{READY}}, I_{\text{STAGE}}, I_{\text{HIER}}) \in [0,1]^9
\end{equation}

\textbf{For Practitioners:} See Appendix HHH (METHOD-TOOLKIT), Section ``Quickstart Guide: Using the IE Master Matrix'' for step-by-step intervention design protocol.

\textbf{Cross-References:} All 10C COREs (each maps to one $I_c$ dimension), HHH (operational tools + Quickstart Guide), Chapters 17-22 (intervention design and policy applications)
\end{tcolorbox}

\subsection{CORE-PRINCIPLES: The Seven Fundamental Axioms (Appendix PR)}

\begin{tcolorbox}[colback=gray!5!white, colframe=gray!75!black]
\textbf{Question:} What are the defining characteristics that make a framework ``EBF-compliant''?

\textbf{Answer:} A framework belongs to the EBF class if and only if it satisfies all seven fundamental principles P1-P7: Complementarity (multi-factor interactions), Context (8-dimensional $\Psi$ specification), Systematik (10C structured analysis), Evidence (scientific grounding), Falsifiability (testable predictions with CI), Human-Machine (Menschlichkeit $\times$ Mechanik amplification), and Cumulation ($O(N \log N)$ knowledge scalability).

\textbf{Key Content:}
\begin{itemize}[nosep]
\item 7 formal axioms (PR-1 to PR-7) defining EBF membership
\item Interdependency matrix showing principle interactions
\item EBF Completeness Theorem: $\mathcal{F} \in \text{EBF} \Leftrightarrow \bigwedge_{i=1}^{7} P_i$
\item Meta-principle: The three amplification layers (Behavioral Theory $\times$ Scientific Method $\times$ Scalability)
\item Worked Example: Swiss retirement savings analyzed through all 7 principles
\end{itemize}

\textbf{Key Formula:}
\begin{equation}
\text{EBF} = \underbrace{(P_1 + P_2 + P_3)}_{\text{Behavioral Theory}} \times \underbrace{(P_4 + P_5)}_{\text{Scientific Method}} \times \underbrace{(P_6 + P_7)}_{\text{Scalability}}
\end{equation}

\textbf{Cross-References:} All 10C COREs (P1-P3 formalize their structure), HI (P6 Human-Machine principle), BX (BEATRIX architecture), FM (scientific reusability), S (P5 falsifiable predictions)
\end{tcolorbox}

\subsection{The 10C CORE Integration: From Utility to Action}

\begin{tcolorbox}[colback=blue!5!white, colframe=blue!75!black,
    title=The Complete EBF Behavioral Framework]

Combining all eight CORE appendices (10C) yields the complete specification:

\textbf{Stage 1: Utility Definition (COREs 1-5)}
\begin{equation}
U^{pot}_{total}(\Psi, \Theta) = \sum_{L=0}^{\infty} \alpha^L(\Psi) \cdot \left[
    \sum_{d \in FEPSDE} \omega^L_d \cdot U^L_d +
    \sum_{d < d'} \gamma^L_{dd'} \cdot U^L_d \cdot U^L_{d'}
\right]
\end{equation}

\textbf{Stage 2: Awareness Filter (CORE 6: AWARE)}
\begin{equation}
U^{eff}_{total}(t^*) = A(t^*) \times U^{pot}_{total}
\end{equation}

\textbf{Stage 3: Action Threshold (CORE 7: READY)}
\begin{equation}
\boxed{\text{Action} \Leftrightarrow WAX(U^{eff}, \varphi, \Psi) \geq \theta(\Psi)}
\end{equation}

\textbf{Stage 4: Journey Phase (CORE 8: STAGE)}
\begin{equation}
\varphi = \text{BCJ\_Phase}(t) \in \{1,...,5\}
\textbf{Stage 4: Behavioral Change Journey (CORE 8)}
\begin{equation}
\frac{dS}{dt} = \frac{1}{\tau} \cdot \left[ S^*(A, WAX, \Psi) - S(t) \right]
\end{equation}

\begin{center}
\begin{tabular}{ll}
$L$ & from CORE-WHO (AAA): Levels \\
$d$ & from CORE-WHAT (C): Dimensions \\
$\gamma$ & from CORE-HOW (B): Interactions \\
$\Psi, \alpha$ & from CORE-WHEN (V): Context \\
$\Theta = \{\omega, \gamma, \alpha, \lambda, \delta\}$ & from CORE-WHERE (BBB): Parameters \\
$A(t^*)$ & from CORE-AWARE (AU): Salience filter \\
$WAX, \theta$ & from CORE-READY (AV): Action threshold \\
$\varphi$ & from CORE-STAGE (AW): BCJ phase \\
$WAX, \theta, \varphi$ & from CORE-READY (AV): Action threshold \\
$S(t), dS/dt$ & from CORE-STAGE (AW): Change journey \\
\end{tabular}
\end{center}

\end{tcolorbox}

%%%%%%%%%%%%%%%%%%%%%%%%%%%%%%%%%%%%%%%%%%%%%%%%%%%%%%%%%%%%%%%%%%%%%%%%%%%%%%%
%%%%%%%%%%%%%%%%%%%%%%%%%%%%%%%%%%%%%%%%%%%%%%%%%%%%%%%%%%%%%%%%%%%%%%%%%%%%%%%
\section{FORMAL: Mathematical Foundations}
\label{sec:index-formal}
%%%%%%%%%%%%%%%%%%%%%%%%%%%%%%%%%%%%%%%%%%%%%%%%%%%%%%%%%%%%%%%%%%%%%%%%%%%%%%%
%%%%%%%%%%%%%%%%%%%%%%%%%%%%%%%%%%%%%%%%%%%%%%%%%%%%%%%%%%%%%%%%%%%%%%%%%%%%%%%

\begin{table}[htbp]
\centering
\caption{FORMAL Appendices: Mathematical Foundations}
\label{tab:index-formal}
\renewcommand{\arraystretch}{1.4}
\begin{tabular}{@{}clll@{}}
\toprule
\textbf{Code} & \textbf{Name} & \textbf{Title} & \textbf{Content} \\
\midrule
A & FORMAL-DERIVE & Mathematical Foundations & Derivations of core equations \\
D & FORMAL-PROOF & Theorem Demonstrations & Formal proofs of propositions \\
BA & FORMAL-SEGMENT & Axioms of Behavioral Change Segments & Formalization of segment heterogeneity \\
BB & FORMAL-EQUILIBRIA & Intervention Equilibria and Tipping Points & Counterfactual scenarios, $(\xi, \kappa_{\text{crit}})$ operationalization \\
BC & FORMAL-PROBABILITY & The 10-Axiom Probabilistic Foundation & Axioms P1-P10 for behavior transition probabilities \\
III & FORMAL-EIH & Efficient Intervention Hypothesis & Efficiency conditions $\eta = \mu \cdot \sigma_s \cdot \prod(1+\gamma)$ \\
IV & FORMAL-LET & Level Emergence Theorem & $\mathcal{E} = U^L - \sum U^{ACE-1}$ when $\bar{\gamma} > 0$ \\
V & FORMAL-MEP & Minimum Effective Portfolio & Markowitz analogue for intervention optimization \\
VI & FORMAL-MEQ & Multiple Equilibria Theorem & Strategic navigation across equilibrium landscape \\
VII & FORMAL-GTC & General Theory of Complementarity & Unified $\gamma$ framework, consistency theorems \\
X & FORMAL-FOUND & Mathematical Foundations of Complementarity & Domain-independent theory, lattice structure, research landscape \\
XI & FORMAL-MODELS & Canonical Complementarity Models & Six mathematical models (A-F), equivalence theorems \\
NC & FORMAL-NC & Non-Concavity, Fit, and Coherence & $\gamma > 0 \Rightarrow$ non-concave landscape, $K$ vs $Q$ distinction \\
CJ & FORMAL-BELIEFARCHITECTURE & Sabatier ACF applied to social democracy & Deep/Policy/Secondary belief hierarchy; Austrian incoherence analysis \\
CK & FORMAL-BELIEFOPPOSITIONS & Ideological contrasts & Social Democratic vs. Liberal vs. Conservative beliefs; coalition feasibility \\
CM & FORMAL-FAIRNESSALIGNMENT & Developmental fairness and voter incoherence detection & Fehr archetypes (Equality, Equity, Need-based) mapped to life-course; fairness-based detection mechanism \\
UN & FORMAL-UNTCM & Unified Natural Trajectory Model & Coupled ODE system (A5, T6-T8), methodology (M1-M4), integrations (I1-I4), extensions (X1-X4) \\
PC & FORMAL-PCT & Parameter Context Transformation & $\theta = F(\Psi, 10C)$, Axioms PCT-1 to PCT-5, iterative Bayesian learning, measurement contexts \\
\bottomrule
\end{tabular}
\end{table}

\begin{tcolorbox}[colback=yellow!5!white, colframe=yellow!75!black]
\textbf{Note:} AU (Awareness) and AV (Willingness) were elevated to CORE status in January 2026.
\begin{itemize}[nosep]
    \item AU $\to$ CORE-AWARE (6th CORE): The Awareness Function
    \item AV $\to$ CORE-READY (7th CORE): The Willingness Function
    \item AW $\to$ CORE-STAGE (8th CORE): The Behavioral Change Journey
\end{itemize}
See Section~\ref{sec:index-core} for details.
\end{tcolorbox}

%%%%%%%%%%%%%%%%%%%%%%%%%%%%%%%%%%%%%%%%%%%%%%%%%%%%%%%%%%%%%%%%%%%%%%%%%%%%%%%
%%%%%%%%%%%%%%%%%%%%%%%%%%%%%%%%%%%%%%%%%%%%%%%%%%%%%%%%%%%%%%%%%%%%%%%%%%%%%%%
\section{DOMAIN: Application Fields}
\label{sec:index-domain}
%%%%%%%%%%%%%%%%%%%%%%%%%%%%%%%%%%%%%%%%%%%%%%%%%%%%%%%%%%%%%%%%%%%%%%%%%%%%%%%
%%%%%%%%%%%%%%%%%%%%%%%%%%%%%%%%%%%%%%%%%%%%%%%%%%%%%%%%%%%%%%%%%%%%%%%%%%%%%%%

\begin{table}[htbp]
\centering
\caption{DOMAIN Appendices: Economics Subfield Applications}
\label{tab:index-domain}
\renewcommand{\arraystretch}{1.3}
\small
\begin{tabular}{@{}cllp{5cm}@{}}
\toprule
\textbf{Code} & \textbf{Name} & \textbf{Field} & \textbf{EBF Application} \\
\midrule
AA & DOMAIN-LABOR & Labor Economics & Multi-level labor market analysis \\
AB & DOMAIN-MATCH & Matching Theory & Market design with welfare hierarchy \\
AC & DOMAIN-IO & Industrial Organization & Firm-level complementarities \\
AD & DOMAIN-EVO & Evolutionary Game Theory & Level selection dynamics \\
AE & DOMAIN-MECH & Mechanism Design & Level-aware incentive design \\
AF & DOMAIN-CHOICE & Social Choice & Aggregation across levels \\
AG & DOMAIN-COMPLEX & Complexity Economics & Emergent level properties \\
AH & DOMAIN-CONSULTING & Behavioral Consulting & Framework vs Model, 20 KPIs \\
AJ & DOMAIN-SOCIAL & Social Preferences & Other-regarding at multiple levels \\
AK & DOMAIN-EPISTEMIC & Epistemics & Belief formation across levels \\
W & DOMAIN-INFO & Information Economics & Information asymmetry by level \\
X & DOMAIN-COMPLEMENT & Milgrom-Roberts & Supermodularity foundations \\
Y & DOMAIN-CAPITAL & Capital Markets & Financial level interactions \\
Z & DOMAIN-GROWTH & Growth Theory & National welfare dynamics \\
BF & DOMAIN-SOCIALDEMOCRACY & Welfare States \& Political Economy & Redistribution preferences, institutions, inequality \\
XII & DOMAIN-FOUND & Domain Foundations of Complementarity & Economics, Psychology, Management translation \\
XIII & DOMAIN-EBF & EBF Complementarity Application & Behavioral interventions, FEPSDE, BCJ \\
XIV & DOMAIN-CATALOG & Universal Domain Map & 20 domains across 4 science clusters \\
BI & DOMAIN-LIFESPAN & Life Course \& Intergenerational & Intergenerational spillover, multi-generational planning \\
CP & DOMAIN-CPS & Complex Problem Solving & Dörner CPS paradigm as 10C configuration \\
BL & DOMAIN-CRISIS & Geopolitical Crisis Escalation & Sequential game theory for geopolitical decisions \\
\bottomrule
\end{tabular}
\end{table}

%%%%%%%%%%%%%%%%%%%%%%%%%%%%%%%%%%%%%%%%%%%%%%%%%%%%%%%%%%%%%%%%%%%%%%%%%%%%%%%
%%%%%%%%%%%%%%%%%%%%%%%%%%%%%%%%%%%%%%%%%%%%%%%%%%%%%%%%%%%%%%%%%%%%%%%%%%%%%%%
\section{CONTEXT: The $\Psi$ Components}
\label{sec:index-context}
%%%%%%%%%%%%%%%%%%%%%%%%%%%%%%%%%%%%%%%%%%%%%%%%%%%%%%%%%%%%%%%%%%%%%%%%%%%%%%%
%%%%%%%%%%%%%%%%%%%%%%%%%%%%%%%%%%%%%%%%%%%%%%%%%%%%%%%%%%%%%%%%%%%%%%%%%%%%%%%

\begin{table}[htbp]
\centering
\caption{CONTEXT Appendices: Detailed $\Psi$ Dimension Analysis}
\label{tab:index-context}
\renewcommand{\arraystretch}{1.4}
\begin{tabular}{@{}clll@{}}
\toprule
\textbf{Code} & \textbf{Name} & \textbf{Title} & \textbf{$\Psi$ Dimension} \\
\midrule
AH & CONTEXT-TIME & Temporal Context & $\Psi_T$: Time pressure, horizons \\
AI & CONTEXT-SPACE & Spatial Context & $\Psi_{Sp}$: Location, distance \\
BG & CONTEXT-GENESIS & Why Austrian Media Concentration Became Possible & Historical-comparative analysis of $\Psi_{\text{media}}$ emergence \\
BH & CONTEXT-SPÖ_PARADOX & Austrian Social Democrats' Ambivalent Role & SPÖ's inaction on media concentration (1950-2025) \\
DX & CONTEXT-MEDIA & Political Economy of Media & $\Psi_{\text{media}}$: Government ads, ownership concentration \\
V & CONTEXT-MASTER & The Context Framework & All 8 $\Psi$ dimensions \\
\bottomrule
\end{tabular}
\end{table}

%%%%%%%%%%%%%%%%%%%%%%%%%%%%%%%%%%%%%%%%%%%%%%%%%%%%%%%%%%%%%%%%%%%%%%%%%%%%%%%
%%%%%%%%%%%%%%%%%%%%%%%%%%%%%%%%%%%%%%%%%%%%%%%%%%%%%%%%%%%%%%%%%%%%%%%%%%%%%%%
\section{METHOD: Estimation and Validation}
\label{sec:index-method}
%%%%%%%%%%%%%%%%%%%%%%%%%%%%%%%%%%%%%%%%%%%%%%%%%%%%%%%%%%%%%%%%%%%%%%%%%%%%%%%
%%%%%%%%%%%%%%%%%%%%%%%%%%%%%%%%%%%%%%%%%%%%%%%%%%%%%%%%%%%%%%%%%%%%%%%%%%%%%%%

\begin{table}[htbp]
\centering
\caption{METHOD Appendices: Methodological Tools}
\label{tab:index-method}
\renewcommand{\arraystretch}{1.4}
\begin{tabular}{@{}cllp{5cm}@{}}
\toprule
\textbf{Code} & \textbf{Name} & \textbf{Title} & \textbf{Method} \\
\midrule
AN & METHOD-LLMMC & LLM Monte Carlo & Parameter estimation via LLM simulation \\
AL & METHOD-SRL & Self-Reinforcement & Learning dynamics modeling \\
E & METHOD-OPS & Operationalization & Measurement protocols \\
R & METHOD-EVAL & Evaluation Protocol & Validation framework \\
AZ & METHOD-CONSTRUCT & Model Construction & Microfoundations, dynamics, identification \\
CCC & METHOD-DOCTYPE & 8D Document Clustering & Document type characterization \\
EEE & METHOD-DESIGN & Model Design Workflow & Systematic model creation process \\
FFF & METHOD-REGISTRY & Model Archive & Organizational learning repository \\
GGG & METHOD-CONFIG & Model Configurator & Dimension-based model generation \\
HHH & METHOD-TOOLKIT & Intervention Design Toolkit & Systematic intervention specification \\
VIII & METHOD-COMP & Complementarity Methodology & 120+ papers on $\gamma$ measurement \\
IX & METHOD-RESEARCH & Research Agenda & 6 open questions, future directions \\
III & METHOD-LLMMC-CAL & Calibration Pipeline & LLMMC $\rightarrow$ R-Score end-to-end \\
XX & METHOD-EVIDENCE & Disciplinary Boundaries & BE vs BS vs NE evidence classification \\
JJJ & METHOD-PSG & Parameter Sourcing Guide & Sourcing, deriving, and documenting parameters in examples \\
KKK & METHOD-DYNAMIC & Portfolio Evolution Dynamics & Formal models for dynamic portfolio redesign: ODEs, Markov, Bayesian learning, convergence \\
LLL & METHOD-POLICY-DOMAINS & Domain-Specific Portfolio Design & Templates for Health, Pension, Environment, Labor domains with pre-filled matrices \\
MMM & METHOD-POLICY-COHERENCE & Cross-Domain Spillovers & Spillover analysis framework, conflict detection, coherence scoring \\
NNN & METHOD-MULTI-LEVEL & Multi-Level Implementation & Templates and parameters for all 5 levels: Individual, Household, Organization, National, Global \\
RRR & METHOD-LIFETIME & Life Journey Domain Templates & Domain-specific 30-year planning templates, diagnostics, KPIs, segment transitions per domain \\
BJ & METHOD-DOMAINVAL & Life Domain Validation & Cross-cultural validation: 20 traditional societies, 15 WEIRD subcultures, Japan/Gulf/Afghanistan applications; N$_{\text{eff}}$ = 2--7 spectrum; PRO/CONTRA literature (35+ papers) \\
VC & METHOD-VALIDATE & Data Consistency Validation & Referential integrity, parameter consistency, context validation; DCV-1 to DCV-6 axioms \\
\midrule
\textbf{FORMAL APPENDICES (Mathematical Formalization)} \\
\midrule
OOO & FORMAL-SYSTEM-DYNAMICS & Complete ODE Systems & Coupled differential equations for all 7 domains with equilibrium analysis, convergence proofs, bifurcation theory \\
\midrule
\textbf{CONTINUED METHOD APPENDICES} \\
\midrule
PPP & METHOD-AGENT-BASED & Agent-Based Modeling & ABM framework with agent architecture, micro-level dynamics, population emergence, policy scenario testing \\
BK & METHOD-CALC & EBF Calculation Reference & Consolidated formulas: P$_{eff}$, Awareness ODEs, Willingness, Intergenerational spillover, N$_{eff}$ scaling, ROI multipliers, Portfolio optimization \\
BM & METHOD-PAPERINT & Paper Integration Methodology & 2D Classification (Content 0--3, Integration 1--5), Learning Loop, upgrade workflow \\
VB & METHOD-PAPERQA & Paper Integration QA & 13-component validation, compliance thresholds, 6-Factor Framework for CORE decisions \\
\bottomrule
\end{tabular}
\end{table}

%%%%%%%%%%%%%%%%%%%%%%%%%%%%%%%%%%%%%%%%%%%%%%%%%%%%%%%%%%%%%%%%%%%%%%%%%%%%%%%
%%%%%%%%%%%%%%%%%%%%%%%%%%%%%%%%%%%%%%%%%%%%%%%%%%%%%%%%%%%%%%%%%%%%%%%%%%%%%%%
\section{PREDICT: Falsifiable Predictions}
\label{sec:index-predict}
%%%%%%%%%%%%%%%%%%%%%%%%%%%%%%%%%%%%%%%%%%%%%%%%%%%%%%%%%%%%%%%%%%%%%%%%%%%%%%%
%%%%%%%%%%%%%%%%%%%%%%%%%%%%%%%%%%%%%%%%%%%%%%%%%%%%%%%%%%%%%%%%%%%%%%%%%%%%%%%

\begin{table}[htbp]
\centering
\caption{PREDICT Appendices: Testable Hypotheses}
\label{tab:index-predict}
\renewcommand{\arraystretch}{1.3}
\begin{tabular}{@{}cllp{5cm}@{}}
\toprule
\textbf{Code} & \textbf{Name} & \textbf{Title} & \textbf{Prediction Focus} \\
\midrule
S & PREDICT-MASTER & Falsifiable Predictions & 10 core testable hypotheses \\
AO & PREDICT-01 & Trade War Dynamics & Cross-national level conflict \\
AP & PREDICT-02 & Spillover Effects & Cross-level utility transmission \\
AQ & PREDICT-03 & Asymmetric Response & Level-dependent reactions \\
AR & PREDICT-04 & Techlash Dynamics & Technology-identity conflict \\
AS & PREDICT-05 & Identity Activation & $IDN^L$ salience triggers \\
AT & PREDICT-06 & Coherence Trap & Cross-level coherence failure \\
HF & PREDICT-HABIT & Habit Formation Mechanisms & Experience-good vs addiction taxonomy \\
\bottomrule
\end{tabular}
\end{table}

%%%%%%%%%%%%%%%%%%%%%%%%%%%%%%%%%%%%%%%%%%%%%%%%%%%%%%%%%%%%%%%%%%%%%%%%%%%%%%%
%%%%%%%%%%%%%%%%%%%%%%%%%%%%%%%%%%%%%%%%%%%%%%%%%%%%%%%%%%%%%%%%%%%%%%%%%%%%%%%
\section{LIT: Literature Integration}
\label{sec:index-lit}
%%%%%%%%%%%%%%%%%%%%%%%%%%%%%%%%%%%%%%%%%%%%%%%%%%%%%%%%%%%%%%%%%%%%%%%%%%%%%%%
%%%%%%%%%%%%%%%%%%%%%%%%%%%%%%%%%%%%%%%%%%%%%%%%%%%%%%%%%%%%%%%%%%%%%%%%%%%%%%%

\begin{table}[htbp]
\centering
\caption{LIT Appendices: Research Integration by Author}
\label{tab:index-lit}
\renewcommand{\arraystretch}{1.3}
\begin{tabular}{@{}cllp{5cm}@{}}
\toprule
\textbf{Code} & \textbf{Name} & \textbf{Researcher(s)} & \textbf{Contribution to EBF} \\
\midrule
I & LIT-NOBEL & Nobel Laureates & Foundational economics contributions \\
J & LIT-RECENT & Recent Papers & 2020--2025 research integration \\
K & LIT-FEHR & Ernst Fehr & Social preferences, cooperation \\
L & LIT-ACEMOGLU & Daron Acemoglu & Institutions, political economy \\
M & LIT-SHLEIFER & Andrei Shleifer & Behavioral finance, law \\
N & LIT-HECKMAN & James Heckman & Human capital, selection \\
O & LIT-AUTOR & David Autor & Labor markets, technology \\
P & LIT-DUFLO & Esther Duflo & Development, RCTs \\
Q & LIT-BLOOM & Nick Bloom & Uncertainty, management \\
R & LIT-THALER & Richard Thaler & Mental accounting, choice architecture, behavioral finance \\
S & LIT-SUNSTEIN & Cass Sunstein & Behavioral law, policy, choice architecture, group behavior \\
T & LIT-CAMERER & Colin Camerer & Neuroeconomics, game theory, strategic thinking \\
W & LIT-ARIELY & Dan Ariely & Irrationality, dishonesty, behavioral interventions \\
X & LIT-LOEWENSTEIN & George Loewenstein & Emotions, visceral influences, curiosity, temporal discounting \\
Y & LIT-CIALDINI & Robert Cialdini & Influence, social proof, persuasion, compliance \\
Z & LIT-HAIDT & Jonathan Haidt & Moral psychology, polarization, group behavior \\
AA & LIT-MULLAINATHAN & Sendhil Mullainathan & Scarcity, poverty, behavioral policy \\
AB & LIT-LIST & John List & Field experiments, market behavior, allocation \\
AC & LIT-DOLAN & Paul Dolan & Wellbeing economics, happiness, experience design \\
AD & LIT-SPECIALISTS & Multiple Specialists & Behavioral finance, game theory, cross-cultural \\
AE & LIT-FALK & Armin Falk & Cooperation, reciprocity, fairness, labor economics \\
AG & LIT-SHAFIR & Eldar Shafir & Decision-making, bounded rationality, poverty, behavioral policy \\
AM & LIT-GAECHTER & Simon Gächter & Reciprocity, cooperation, punishment, social norms, experimental economics \\
BD & LIT-ALESINA & Alberto Alesina & Redistribution preferences, ethnic heterogeneity, political institutions, social democracy \\
SU & LIT-SUTTER & Matthias Sutter & Team decision making, developmental economics, workplace interventions \\
CA & LIT-CAMERER & Colin Camerer & Behavioral game theory, cognitive hierarchy, neuroeconomics \\
LO & LIT-LOEWENSTEIN & George Loewenstein & Visceral influences, hot-cold empathy gap, curiosity, emotions \\
GI & LIT-GIGERENZER & Gerd Gigerenzer & Fast-and-frugal heuristics, ecological rationality, risk literacy \\
RO & LIT-ROTH & Alvin Roth & Market design, matching theory, kidney exchange, repugnance \\
RA & LIT-RABIN & Matthew Rabin & Reference-dependent preferences, present bias, belief biases \\
FI & LIT-FISCHBACHER & Urs Fischbacher & Conditional cooperation, z-Tree, experimental methods \\
MM & LIT-MARECHAL & Michel Maréchal & Honesty, corporate culture, identity priming, civic behavior \\
LA & LIT-LAIBSON & David Laibson & Hyperbolic discounting, present bias, defaults, commitment \\
GN & LIT-GNEEZY & Uri Gneezy & Incentives, crowding out, gender competition, deception \\
DV & LIT-DELLAVIGNA & Stefano DellaVigna & Contracts, media effects, field experiments, scale attenuation \\
SA & LIT-AMBÜHL & Sandro Ambühl & Beliefs, organ markets, informed consent, paternalism \\
MN & LIT-NIEDERLE & Muriel Niederle & Gender competition, affirmative action, market design \\
XG & LIT-GABAIX & Xavier Gabaix & Behavioral inattention, sparsity, power laws \\
BE & LIT-ENKE & Benjamin Enke & Moral universalism, cognitive uncertainty, culture \\
PA & LIT-AGHION & Philippe Aghion & Schumpeterian growth, creative destruction, innovation \\
DC & LIT-CARD & David Card & Labor economics, natural experiments, minimum wage \\
GA & LIT-AKERLOF & George Akerlof & Identity economics, asymmetric information, behavioral macro \\
VS & LIT-SMITH & Vernon Smith & Experimental economics, market design, two rationalities \\
TS & LIT-SCHELLING & Thomas Schelling & Focal points, segregation, self-command, commitment \\
GB & LIT-BECKER & Gary Becker & Human capital, discrimination, crime, family, rational addiction \\
JT & LIT-TIROLE & Jean Tirole & Industrial organization, platforms, regulation, behavioral IO \\
RB & LIT-BENABOU & Roland Bénabou & Motivation, identity, inequality, groupthink \\
SE & LIT-STRATIFICATION & Darity, Hamilton et al. & Intergroup inequality, wealth gaps, Baby Bonds, reparations \\
OT & LIT-OTHER & Miscellaneous & Catch-all for papers without dedicated author appendix \\
EV & LIT-VANDENSTEEN & Eric Van den Steen & Corporate culture, vision, strategy, organizational economics \\
RS & LIT-SADUN & Raffaella Sadun & Management practices, family firms, healthcare, hybrid work \\
CG & LIT-GOLDIN & Claudia Goldin & Gender economics, labor history, education, greedy jobs (Nobel 2023) \\
RN & LIT-NAGEL & Rosemarie Nagel & Level-k reasoning, beauty contests, cognitive hierarchy, bounded rationality \\

U & LIT-KT & Kahneman-Tversky & Prospect theory foundations \\
AX & LIT-META & Metascience & Why integration is needed \\
AY & LIT-PARADIGMS & Multi-disciplinary & Reference points across frameworks \\
XV & LIT-HISTORY & Scientific History & 2,500-year trajectory of complementarity \\
XVI & LIT-CRITIQUE & Critical Literature & 10 canonical arguments and 100 references \\
XVII & LIT-THERMODYNAMICS & Economic Impact & 5 functions, 25 cases, $\gamma$-analogy \\
XVIII & LIT-FEHR-METHOD & Integration vs. Falsifiability & Fehr's methodological critique, $\gamma = 0$ null hypothesis, disciplines all 10C \\
XIX & LIT-LEARNING & Experience Goods \& Belief Updating & Nelson, Ackerberg, Jones et al., learning vs. addiction \\
\bottomrule
\end{tabular}
\end{table}

%%%%%%%%%%%%%%%%%%%%%%%%%%%%%%%%%%%%%%%%%%%%%%%%%%%%%%%%%%%%%%%%%%%%%%%%%%%%%%%
%%%%%%%%%%%%%%%%%%%%%%%%%%%%%%%%%%%%%%%%%%%%%%%%%%%%%%%%%%%%%%%%%%%%%%%%%%%%%%%
\section{REF: Reference Materials}
\label{sec:index-ref}
%%%%%%%%%%%%%%%%%%%%%%%%%%%%%%%%%%%%%%%%%%%%%%%%%%%%%%%%%%%%%%%%%%%%%%%%%%%%%%%
%%%%%%%%%%%%%%%%%%%%%%%%%%%%%%%%%%%%%%%%%%%%%%%%%%%%%%%%%%%%%%%%%%%%%%%%%%%%%%%

\begin{table}[htbp]
\centering
\caption{REF Appendices: Reference Documentation}
\label{tab:index-ref}
\renewcommand{\arraystretch}{1.4}
\begin{tabular}{@{}clll@{}}
\toprule
\textbf{Code} & \textbf{Name} & \textbf{Title} & \textbf{Content} \\
\midrule
F & REF-EXAMPLES & Worked Examples & Step-by-step applications \\
G & REF-GLOSSARY & Complete Glossary & All terms and symbols \\
H & REF-HISTORY & Computational History & Development timeline \\
T & REF-META & Metatheory & Philosophical foundations \\
SK & REF-SKILLS & Skill Registry & Claude Code skills with EBF integration \\
DDD & REF-DOCTYPE & Document Type Theory & 8D audience-centered framework \\
DA & REF-DEFENSE & Critical Questions \& Defense & 20 killer questions (P1-P10, W1-W10) \\
DA2 & REF-DEFENSE-MATRIX & Response Variants & 6 variants per question ($3 \times 2$) \\
BO & REF-ARCH & Component Architecture & EBF Gefässe \& SSOT Registry \\
BN & REF-LEAD & Lead Management System & 6-stage lifecycle, 3-layer ID architecture, YAML data model, workflows \\
\bottomrule
\end{tabular}
\end{table}

%%%%%%%%%%%%%%%%%%%%%%%%%%%%%%%%%%%%%%%%%%%%%%%%%%%%%%%%%%%%%%%%%%%%%%%%%%%%%%%
%%%%%%%%%%%%%%%%%%%%%%%%%%%%%%%%%%%%%%%%%%%%%%%%%%%%%%%%%%%%%%%%%%%%%%%%%%%%%%%
\section{Complete Appendix Map}
\label{sec:index-complete}
%%%%%%%%%%%%%%%%%%%%%%%%%%%%%%%%%%%%%%%%%%%%%%%%%%%%%%%%%%%%%%%%%%%%%%%%%%%%%%%
%%%%%%%%%%%%%%%%%%%%%%%%%%%%%%%%%%%%%%%%%%%%%%%%%%%%%%%%%%%%%%%%%%%%%%%%%%%%%%%

\begin{table}[htbp]
\centering
\caption{Complete Appendix Index: All 78 Appendices (with 10 CORE)}
\label{tab:index-complete}
\renewcommand{\arraystretch}{1.1}
\scriptsize
\begin{tabular}{@{}cllc@{}}
\toprule
\textbf{Code} & \textbf{Full Name} & \textbf{Category} & \textbf{Priority} \\
\midrule
\multicolumn{4}{@{}l}{\textit{\textbf{CORE (9)}}} \\
AAA & CORE-WHO: The Welfare Hierarchy & Core Theory & Essential \\
C & CORE-WHAT: The Welfare Dimensions & Core Theory & Essential \\
B & CORE-HOW: The Interaction Structure & Core Theory & Essential \\
V & CORE-WHEN: The Context Framework & Core Theory & Essential \\
BBB & CORE-WHERE: Parameter Estimation & Core Theory & Essential \\
AU & CORE-AWARE: The Awareness Function & Core Theory & Essential \\
AV & CORE-READY: The Willingness Function & Core Theory & Essential \\
AW & CORE-STAGE: The Behavioral Change Journey & Core Theory & Essential \\
HI & CORE-HIERARCHY: Decision Hierarchy \& Strategic Coherence & Core Theory & Essential \\
PR & CORE-PRINCIPLES: Fundamental Axioms of EBF & Core Theory & Essential \\
\addlinespace
\multicolumn{4}{@{}l}{\textit{\textbf{FORMAL (12)}}} \\
A & FORMAL-DERIVE: Mathematical Foundations & Formalization & High \\
D & FORMAL-PROOF: Theorem Demonstrations & Formalization & High \\
BA & FORMAL-SEGMENT: Axioms of Behavioral Change Segments & Formalization & High \\
BB & FORMAL-EQUILIBRIA: Intervention Equilibria \& Tipping Points & Formalization & Essential \\
III & FORMAL-EIH: Efficient Intervention Hypothesis & Formalization & Essential \\
IV & FORMAL-LET: Level Emergence Theorem & Formalization & Essential \\
V & FORMAL-MEP: Minimum Effective Portfolio & Formalization & Essential \\
VI & FORMAL-MEQ: Multiple Equilibria \& Strategic Navigation & Formalization & Essential \\
VII & FORMAL-GTC: General Theory of Complementarity & Formalization & Essential \\
X & FORMAL-FOUND: Mathematical Foundations of Complementarity & Formalization & Essential \\
XI & FORMAL-MODELS: Canonical Complementarity Models & Formalization & Essential \\
NC & FORMAL-NC: Non-Concavity, Fit, and Coherence & Formalization & High \\
UN & FORMAL-UNTCM: Unified Natural Trajectory Model & Formalization & High \\
PC & FORMAL-PCT: Parameter Context Transformation & Formalization & Essential \\
\addlinespace
\multicolumn{4}{@{}l}{\textit{\textbf{DOMAIN (17)}}} \\
AA & DOMAIN-LABOR: Labor Market Applications & Domain & Medium \\
AB & DOMAIN-MATCH: Matching \& Market Design & Domain & Medium \\
AC & DOMAIN-IO: Industrial Organization & Domain & Medium \\
AD & DOMAIN-EVO: Evolutionary Game Theory & Domain & Medium \\
AE & DOMAIN-MECH: Mechanism Design & Domain & Medium \\
AF & DOMAIN-CHOICE: Social Choice Theory & Domain & Medium \\
AG & DOMAIN-COMPLEX: Complexity Economics & Domain & Medium \\
AH & DOMAIN-CONSULTING: Behavioral Consulting & Domain & High \\
AJ & DOMAIN-SOCIAL: Social Preferences & Domain & High \\
AK & DOMAIN-EPISTEMIC: Epistemics \& Beliefs & Domain & Medium \\
W & DOMAIN-INFO: Information Economics & Domain & Medium \\
X & DOMAIN-COMPLEMENT: Milgrom-Roberts & Domain & High \\
Y & DOMAIN-CAPITAL: Capital Markets & Domain & Medium \\
Z & DOMAIN-GROWTH: Growth Theory & Domain & Medium \\
XII & DOMAIN-FOUND: Domain Foundations of Complementarity & Domain & Essential \\
XIII & DOMAIN-EBF: EBF Complementarity Application & Domain & Essential \\
XIV & DOMAIN-CATALOG: Universal Domain Map of Complementarity & Domain & High \\
BI & DOMAIN-LIFESPAN: Life Course \& Intergenerational Effects & Domain & High \\
\addlinespace
\multicolumn{4}{@{}l}{\textit{\textbf{CONTEXT (6)}}} \\
AH & CONTEXT-TIME: Temporal Context & Context & Medium \\
AI & CONTEXT-SPACE: Spatial Context & Context & Medium \\
BG & CONTEXT-GENESIS: Why Austrian Media Concentration Became Possible & Context & High \\
BH & CONTEXT-SPÖ_PARADOX: Austrian Social Democrats' Ambivalent Role & Context & High \\
DX & CONTEXT-MEDIA: Political Economy of Media Markets & Context & Essential \\
V & CONTEXT-MASTER (= CORE-WHEN) & Context & Essential \\
\addlinespace
\multicolumn{4}{@{}l}{\textit{\textbf{METHOD (14)}}} \\
AN & METHOD-LLMMC: LLM Monte Carlo & Methodology & High \\
AL & METHOD-SRL: Self-Reinforcement Learning & Methodology & Medium \\
E & METHOD-OPS: Operationalization Protocol & Methodology & High \\
R & METHOD-EVAL: Evaluation Framework & Methodology & High \\
AZ & METHOD-CONSTRUCT: Model Construction Foundations & Methodology & High \\
CCC & METHOD-DOCTYPE: 8D Document Clustering & Methodology & Medium \\
EEE & METHOD-DESIGN: Model Design Workflow & Methodology & High \\
FFF & METHOD-REGISTRY: Model Archive & Methodology & High \\
GGG & METHOD-CONFIG: Model Configurator & Methodology & High \\
HHH & METHOD-TOOLKIT: Intervention Design Toolkit & Methodology & High \\
VIII & METHOD-COMP: Complementarity Methodology & Methodology & Essential \\
IX & METHOD-RESEARCH: Complementarity Research Agenda & Methodology & High \\
III & METHOD-LLMMC-CAL: Calibration Pipeline & Methodology & High \\
XX & METHOD-EVIDENCE: Disciplinary Boundaries & Methodology & High \\
JJJ & METHOD-PSG: Parameter Sourcing Guide & Methodology & High \\
KKK & METHOD-DYNAMIC: Portfolio Evolution Dynamics & Methodology & High \\
LLL & METHOD-POLICY-DOMAINS: Domain-Specific Portfolio Design & Methodology & High \\
MMM & METHOD-POLICY-COHERENCE: Cross-Domain Spillovers & Methodology & High \\
NNN & METHOD-MULTI-LEVEL: Multi-Level Implementation & Methodology & High \\
RRR & METHOD-LIFETIME: Life Journey Domain Templates & Methodology & High \\
BJ & METHOD-DOMAINVAL: Life Domain Validation & Methodology & High \\
OOO & FORMAL-SYSTEM-DYNAMICS: Complete ODE Systems & Formalization & Essential \\
PPP & METHOD-AGENT-BASED: Agent-Based Modeling & Methodology & High \\
BK & METHOD-CALC: EBF Calculation Reference & Methodology & High \\
BM & METHOD-PAPERINT: Paper Integration Methodology & Methodology & High \\
VC & METHOD-VALIDATE: Data Consistency Validation & Methodology & High \\
\addlinespace
\multicolumn{4}{@{}l}{\textit{\textbf{PREDICT (8)}}} \\
S & PREDICT-MASTER: Falsifiable Predictions & Predictions & Essential \\
AO & PREDICT-01: Trade War Dynamics & Predictions & Medium \\
AP & PREDICT-02: Cross-Level Spillovers & Predictions & Medium \\
AQ & PREDICT-03: Asymmetric Response & Predictions & Medium \\
AR & PREDICT-04: Techlash Dynamics & Predictions & Medium \\
AS & PREDICT-05: Identity Activation & Predictions & Medium \\
AT & PREDICT-06: Coherence Trap & Predictions & Medium \\
HF & PREDICT-HABIT: Habit Formation Mechanisms & Predictions & High \\
\addlinespace
\multicolumn{4}{@{}l}{\textit{\textbf{LIT (30)}}} \\
I & LIT-NOBEL: Nobel Contributions & Literature & Medium \\
J & LIT-RECENT: Recent Papers 2020--2025 & Literature & High \\
K & LIT-FEHR: Ernst Fehr Research & Literature & Essential \\
L & LIT-ACEMOGLU: Daron Acemoglu Research & Literature & High \\
M & LIT-SHLEIFER: Andrei Shleifer Research & Literature & Medium \\
N & LIT-HECKMAN: James Heckman Research & Literature & Medium \\
O & LIT-AUTOR: David Autor Research & Literature & Medium \\
P & LIT-DUFLO: Esther Duflo Research & Literature & Medium \\
Q & LIT-BLOOM: Nick Bloom Research & Literature & Medium \\
R & Richard Thaler Research: Mental Accounting \& Choice Architecture & Literature & High \\
S & Cass Sunstein Research: Behavioral Law \& Policy & Literature & High \\
T & Colin Camerer Research: Neuroeconomics \& Game Theory & Literature & High \\
W & Dan Ariely Research: Irrationality \& Decision-Making & Literature & High \\
X & George Loewenstein Research: Emotions \& Visceral Effects & Literature & High \\
Y & Robert Cialdini Research: Influence \& Social Proof & Literature & High \\
Z & Jonathan Haidt Research: Moral Psychology \& Polarization & Literature & High \\
AA & Sendhil Mullainathan Research: Scarcity \& Behavioral Policy & Literature & High \\
AB & John List Research: Field Experiments \& Market Behavior & Literature & High \\
AC & Paul Dolan Research: Wellbeing \& Experience Design & Literature & High \\
AD & Behavioral Specialists: Finance, Game Theory, Cross-Cultural & Literature & Medium \\
AE & Armin Falk Research: Cooperation, Reciprocity, Fairness & Literature & High \\
AF & Ulrike Malmendier Research: Behavioral Finance \& CEO Overconfidence & Literature & High \\
AG & Eldar Shafir Research: Decision-Making \& Bounded Rationality & Literature & High \\
SU & LIT-SUTTER: Team Decision Making \& Developmental Economics & Literature & High \\
CA & LIT-CAMERER: Behavioral Game Theory \& Neuroeconomics & Literature & High \\
LO & LIT-LOEWENSTEIN: Visceral Influences \& Emotions & Literature & High \\
GI & LIT-GIGERENZER: Ecological Rationality \& Heuristics & Literature & High \\
RO & LIT-ROTH: Market Design \& Matching Theory & Literature & High \\
RA & LIT-RABIN: Reference-Dependent Preferences \& Present Bias & Literature & Essential \\
FI & LIT-FISCHBACHER: Conditional Cooperation \& Experimental Methods & Literature & High \\
MM & LIT-MARECHAL: Honesty \& Corporate Culture & Literature & High \\
LA & LIT-LAIBSON: Hyperbolic Discounting \& Present Bias & Literature & Essential \\
GN & LIT-GNEEZY: Incentives, Gender \& Deception & Literature & High \\
DV & LIT-DELLAVIGNA: Contracts, Media \& Field Evidence & Literature & High \\
SA & LIT-AMBÜHL: Beliefs, Organ Markets \& Informed Consent & Literature & High \\
MN & LIT-NIEDERLE: Gender, Competition \& Market Design & Literature & High \\
XG & LIT-GABAIX: Behavioral Inattention \& Sparsity & Literature & High \\
BE & LIT-ENKE: Moral Universalism \& Cognitive Uncertainty & Literature & High \\
PA & LIT-AGHION: Schumpeterian Growth \& Creative Destruction & Literature & Essential \\
DC & LIT-CARD: Labor Economics \& Natural Experiments & Literature & Essential \\
GA & LIT-AKERLOF: Identity Economics \& Asymmetric Information & Literature & Essential \\
VS & LIT-SMITH: Experimental Economics \& Market Design & Literature & Essential \\
TS & LIT-SCHELLING: Focal Points, Segregation \& Self-Command & Literature & Essential \\
GB & LIT-BECKER: Human Capital, Discrimination \& Rational Addiction & Literature & Essential \\
JT & LIT-TIROLE: Industrial Organization, Platforms \& Behavioral IO & Literature & Essential \\
RB & LIT-BENABOU: Motivation, Identity \& Inequality & Literature & Essential \\
SE & LIT-STRATIFICATION: Intergroup Inequality \& Wealth Gaps & Literature & Essential \\
OT & LIT-OTHER: Miscellaneous Research Collection & Literature & Medium \\
EV & LIT-VANDENSTEEN: Corporate Culture \& Organizational Economics & Literature & High \\
RS & LIT-SADUN: Management Practices \& Organizational Economics & Literature & High \\
CG & LIT-GOLDIN: Gender Economics \& Labor History (Nobel 2023) & Literature & Essential \\
RN & LIT-NAGEL: Level-k Reasoning \& Beauty Contests & Literature & High \\

U & LIT-KT: Kahneman-Tversky Foundations & Literature & Essential \\
AX & LIT-META: Metascience \& Integration & Literature & High \\
AY & LIT-PARADIGMS: Reference Points Across Frameworks & Literature & High \\
XV & LIT-HISTORY: Scientific History of Complementarity & Literature & Essential \\
XVI & LIT-CRITIQUE: Critical Literature \& 10 Arguments & Literature & Essential \\
XVII & LIT-THERMODYNAMICS: Economic Impact of Thermodynamics & Literature & High \\
XVIII & LIT-FEHR-METHOD: Integration vs. Falsifiability & Literature & Essential \\
XIX & LIT-LEARNING: Experience Goods \& Belief Updating & Literature & High \\
\addlinespace
\multicolumn{4}{@{}l}{\textit{\textbf{REF (9)}}} \\
F & REF-EXAMPLES: Worked Examples & Reference & High \\
G & REF-GLOSSARY: Complete Glossary & Reference & Essential \\
H & REF-HISTORY: Computational History & Reference & Low \\
T & REF-META: Metatheoretical Foundations & Reference & Medium \\
SK & REF-SKILLS: EBF Skill Registry & Reference & High \\
DDD & REF-DOCTYPE: Document Type Theory & Reference & Medium \\
DA & REF-DEFENSE: Critical Questions \& Framework Defense & Reference & Essential \\
DA2 & REF-DEFENSE-MATRIX: Response Variants Matrix & Reference & Essential \\
BO & REF-ARCH: Component Architecture \& SSOT Registry & Reference & Essential \\
BN & REF-LEAD: Lead Management System & Reference & High \\
\bottomrule
\end{tabular}
\end{table}

%%%%%%%%%%%%%%%%%%%%%%%%%%%%%%%%%%%%%%%%%%%%%%%%%%%%%%%%%%%%%%%%%%%%%%%%%%%%%%%
\section{Reading Pathways}
\label{sec:index-pathways}
%%%%%%%%%%%%%%%%%%%%%%%%%%%%%%%%%%%%%%%%%%%%%%%%%%%%%%%%%%%%%%%%%%%%%%%%%%%%%%%

\begin{tcolorbox}[colback=green!5!white, colframe=green!75!black, 
    title=Recommended Reading Pathways]

\textbf{Pathway 1: Core Theory (Essential)}
\begin{enumerate}[nosep]
\item CORE-WHO (AAA) → Understand levels
\item CORE-WHAT (C) → Understand dimensions
\item CORE-HOW (B) → Understand interactions
\item CORE-WHEN (V) → Understand context
\item CORE-STAGE (AW) → Understand behavioral change journey
\end{enumerate}

\textbf{Pathway 2: Empirical Researcher}
\begin{enumerate}[nosep]
\item CORE-WHO + CORE-WHAT → Basic framework
\item METHOD-OPS (E) → Measurement
\item METHOD-LLMMC (AN) → Estimation
\item PREDICT-MASTER (S) → Testable hypotheses
\end{enumerate}

\textbf{Pathway 3: Domain Specialist}
\begin{enumerate}[nosep]
\item CORE appendices → Foundation
\item Relevant DOMAIN appendix → Your field
\item LIT appendices → Key researchers
\end{enumerate}

\textbf{Pathway 4: Quick Reference}
\begin{enumerate}[nosep]
\item REF-GLOSSARY (G) → Terminology
\item REF-EXAMPLES (F) → Worked applications
\item Specific appendix as needed
\end{enumerate}

\end{tcolorbox}

%%%%%%%%%%%%%%%%%%%%%%%%%%%%%%%%%%%%%%%%%%%%%%%%%%%%%%%%%%%%%%%%%%%%%%%%%%%%%%%
\section{Cross-Reference Matrix}
\label{sec:index-crossref}
%%%%%%%%%%%%%%%%%%%%%%%%%%%%%%%%%%%%%%%%%%%%%%%%%%%%%%%%%%%%%%%%%%%%%%%%%%%%%%%

\begin{tcolorbox}[colback=yellow!5!white, colframe=yellow!75!black, 
    title=Key Cross-References]

\textbf{CORE-WHO (AAA) references:}
\begin{itemize}[nosep]
\item CORE-WHAT (C): Dimensions at each level
\item CORE-HOW (B): Cross-level complementarity
\item CORE-WHEN (V): Level salience function
\item LIT-FEHR (K): Empirical foundations
\end{itemize}

\textbf{CORE-WHAT (C) references:}
\begin{itemize}[nosep]
\item CORE-WHO (AAA): Level-specific dimensions
\item LIT-KT (U): Prospect theory foundations
\item FORMAL-AXIOM (AU): Axiom system
\end{itemize}

\textbf{PREDICT-MASTER (S) references:}
\begin{itemize}[nosep]
\item All CORE appendices: Theoretical basis
\item METHOD-OPS (E): Measurement protocols
\item All PREDICT-01 through PREDICT-06: Specific cases
\end{itemize}

\textbf{FORMAL-UNTCM (UN) references:}
\begin{itemize}[nosep]
\item FORMAL-PROBABILITY (PBB): NTM foundations (A1-A4, Theorems 1-5)
\item CORE-WHEN (V): Context vector $\vec{\Psi}$ definition
\item CORE-STAGE (AW): BCJ phase mapping (I1)
\item METHOD-TOOLKIT (HHH): 10C intervention mapping (I2)
\item CORE-WHAT (C): FEPSDE decomposition (I3)
\end{itemize}

\end{tcolorbox}

%%%%%%%%%%%%%%%%%%%%%%%%%%%%%%%%%%%%%%%%%%%%%%%%%%%%%%%%%%%%%%%%%%%%%%%%%%%%%%%
\section{Version History}
\label{sec:index-version}
%%%%%%%%%%%%%%%%%%%%%%%%%%%%%%%%%%%%%%%%%%%%%%%%%%%%%%%%%%%%%%%%%%%%%%%%%%%%%%%

\begin{table}[htbp]
\centering
\caption{Appendix Index Version History}
\label{tab:index-version}
\renewcommand{\arraystretch}{1.3}
\begin{tabular}{@{}lll@{}}
\toprule
\textbf{Version} & \textbf{Date} & \textbf{Changes} \\
\midrule
1.0 & January 2026 & Initial comprehensive index \\
 & & 49 appendices catalogued \\
 & & 8 categories established \\
 & & Naming convention defined \\
\bottomrule
\end{tabular}
\end{table}

%%%%%%%%%%%%%%%%%%%%%%%%%%%%%%%%%%%%%%%%%%%%%%%%%%%%%%%%%%%%%%%%%%%%%%%%%%%%%%%
%%%%%%%%%%%%%%%%%%%%%%%%%%%%%%%%%%%%%%%%%%%%%%%%%%%%%%%%%%%%%%%%%%%%%%%%%%%%%%%
\section{Chapter-Appendix Mapping}
\label{sec:index-chapter-mapping}
%%%%%%%%%%%%%%%%%%%%%%%%%%%%%%%%%%%%%%%%%%%%%%%%%%%%%%%%%%%%%%%%%%%%%%%%%%%%%%%
%%%%%%%%%%%%%%%%%%%%%%%%%%%%%%%%%%%%%%%%%%%%%%%%%%%%%%%%%%%%%%%%%%%%%%%%%%%%%%%

\begin{tcolorbox}[colback=purple!5!white, colframe=purple!75!black, 
    title=Critical: Every Appendix Links to Main Chapters]

The EBF document follows a dual structure:
\begin{itemize}[nosep]
\item \textbf{Main Chapters} (1--25): Conceptual presentation, narrative flow
\item \textbf{Appendices} (152): Technical details, formal specifications, proofs
\end{itemize}

Every appendix MUST specify:
\begin{enumerate}[nosep]
\item Its \textbf{Primary Chapter} (main conceptual home)
\item Any \textbf{Secondary Chapters} that reference it
\item A \textbf{Section-by-Section mapping} showing correspondence
\end{enumerate}

\end{tcolorbox}

\subsection{Main Document Chapter Structure}

\begin{table}[htbp]
\centering
\caption{EBF Main Document Chapters}
\label{tab:main-chapters}
\renewcommand{\arraystretch}{1.2}
\small
\begin{tabular}{@{}clp{6cm}@{}}
\toprule
\textbf{Ch.} & \textbf{Title} & \textbf{Core Content} \\
\midrule
1 & Introduction & Framework overview, motivation \\
2 & Rationality as Stability & Redefining rationality \\
3 & Limits of Classical Utility & Why extensions needed \\
4 & Empirical Foundations & Evidence base \\
4x & Calibration not Simulation & LLM Monte Carlo method \\
5 & Complementarity & gamma interactions \\
6 & Reference Structure & Reference points \\
7 & Fit and Non-Concavity & Functional forms \\
8 & Mathematical Foundations & Formal framework \\
9 & Context as Endogenous & The Psi framework \\
10 & Die Nutzenarchitektur & INU/KNU/IDN, 144 components, FEPSDE \\
11 & Awareness & Decision processes, A(.) function \\
12 & Willingness & Action thresholds, WAX >= theta \\
13 & Behavioral Change Journey & Stage dynamics, phi phases \\
14 & Behavioral Change Segments & Segment profiles \\
15 & WEC Synthesis & Willingness x Journey x Segment \\
16 & Probability of Behavior Change & Effective probability output \\
17 & Intervention Foundations & Emergent 10C intervention vectors, 20-field schema \\
18 & Journey-Integrated Interventions & Phase-type affinity matrix \\
19 & Segment-Specific Interventions & Segment-type multipliers \\
20 & Intervention Portfolios & Multi-intervention optimization \\
21 & Dynamic Portfolio Evolution & Adaptive rebalancing \\
22 & Policy Applications & Real-world implementation \\
23 & Multi-Level Implementation & Organizational scaling \\
24 & Emergent Life Journeys & Lifespan dynamics, intergenerational spillover \\
\bottomrule
\end{tabular}
\end{table}

\subsection{CORE Appendix $\leftrightarrow$ Chapter Mapping}

\begin{table}[htbp]
\centering
\caption{CORE Appendices: Chapter Linkages}
\label{tab:core-chapter-mapping}
\renewcommand{\arraystretch}{1.3}
\begin{tabular}{@{}llp{7cm}@{}}
\toprule
\textbf{CORE} & \textbf{Primary Chapter} & \textbf{Secondary Chapters} \\
\midrule
\textbf{CORE-WHO} (AAA) & Ch. 10.8: Aggregate Welfare & 
    Ch. 9: Context (level salience) \newline
    Ch. 11: Awareness (level-specific) \\
\addlinespace
\textbf{CORE-WHAT} (C) & Ch. 10: Nutzenarchitektur &
    Ch. 10.1--10.7: All subsections \newline
    Ch. 6: Reference Structure \\
\addlinespace
\textbf{CORE-HOW} (B) & Ch. 5: Complementarity & 
    Ch. 10.6: Inter-Category Compl. \newline
    Ch. 7: Fit and Non-Concavity \\
\addlinespace
\textbf{CORE-WHEN} (V) & Ch. 9: Context & 
    Ch. 10.7: Context Modulation \newline
    Ch. 11.15: $\Psi$-Modulation \\
\bottomrule
\end{tabular}
\end{table}

\subsection{Complete Appendix-Chapter Matrix}

\begin{longtable}{@{}llll@{}}
\caption{All Appendices: Chapter Linkages} 
\label{tab:all-chapter-mapping} \\
\toprule
\textbf{Code} & \textbf{Appendix Name} & \textbf{Primary Ch.} & \textbf{Secondary} \\
\midrule
\endfirsthead
\multicolumn{4}{c}{\tablename\ \thetable{} -- continued} \\
\toprule
\textbf{Code} & \textbf{Appendix Name} & \textbf{Primary Ch.} & \textbf{Secondary} \\
\midrule
\endhead
\midrule
\multicolumn{4}{r}{Continued...} \\
\endfoot
\bottomrule
\endlastfoot

\multicolumn{4}{@{}l}{\textit{\textbf{CORE (8)}}} \\
AAA & CORE-WHO & Ch. 10.8 & Ch. 9, 11 \\
C & CORE-WHAT & Ch. 10 & Ch. 6 \\
B & CORE-HOW & Ch. 5 & Ch. 10.6, 7 \\
V & CORE-WHEN & Ch. 9 & Ch. 10.7, 11.15 \\
BBB & CORE-WHERE & Ch. 4x & Ch. 4, 8 \\
AU & CORE-AWARE & Ch. 11 & Ch. 10, 12 \\
AV & CORE-READY & Ch. 12 & Ch. 11 \\
AW & CORE-STAGE & Ch. 13 & Ch. 14, 15 \\
PR & CORE-PRINCIPLES & Ch. 1 & Ch. 5, 9, HI, BX \\
\addlinespace

\multicolumn{4}{@{}l}{\textit{\textbf{FORMAL (12)}}} \\
A & FORMAL-DERIVE & Ch. 8 & Ch. 5, 10 \\
D & FORMAL-PROOF & Ch. 8 & All \\
BA & FORMAL-SEGMENT & Ch. 14 & Ch. 11, 12, 13 \\
BB & FORMAL-EQUILIBRIA & Ch. 15 & Ch. 14, 13, 12 \\
III & FORMAL-EIH & Ch. 17 & Ch. 15, HHH \\
IV & FORMAL-LET & Ch. 5, 10 & AAA, B, HI \\
V & FORMAL-MEP & Ch. 15, 17 & HHH, III, IV \\
VI & FORMAL-MEQ & Ch. 15, 17 & V, BB, HI, HHH \\
VII & FORMAL-GTC & Ch. 5, 8 & B, All FORMAL \\
X & FORMAL-FOUND & Ch. 8 & VII, VIII, XI \\
XI & FORMAL-MODELS & Ch. 8 & VII, X, B \\
NC & FORMAL-NC & Ch. 7 & B, MIL, VII, IE \\
PC & FORMAL-PCT & Ch. 9 & V, BBB, ALL CORE \\
\addlinespace

\multicolumn{4}{@{}l}{\textit{\textbf{DOMAIN (17)}}} \\
AA & DOMAIN-LABOR & Ch. 10.9.1 & Ch. 4 \\
AB & DOMAIN-MATCH & Ch. 5 & Ch. 10.6 \\
AC & DOMAIN-IO & Ch. 5 & Ch. 7 \\
AD & DOMAIN-EVO & Ch. 3 & Ch. 5 \\
AE & DOMAIN-MECH & Ch. 5 & Ch. 9 \\
AF & DOMAIN-CHOICE & Ch. 3 & Ch. 10.8 \\
AG & DOMAIN-COMPLEX & Ch. 9 & Ch. 5 \\
AH & DOMAIN-CONSULTING & LIT-META & Ch. 14 \\
AJ & DOMAIN-SOCIAL & Ch. 10.1 & Ch. 4 \\
AK & DOMAIN-EPISTEMIC & Ch. 11.03 & Ch. 9 \\
W & DOMAIN-INFO & Ch. 9 & Ch. 5 \\
X & DOMAIN-COMPLEMENT & Ch. 5 & Ch. 7 \\
Y & DOMAIN-CAPITAL & Ch. 10.2 & Ch. 5 \\
Z & DOMAIN-GROWTH & Ch. 10.8 & Ch. 9 \\
XII & DOMAIN-FOUND & Ch. 5, 8 & X, XI, VII \\
XIII & DOMAIN-EBF & Ch. 15, 17 & XII, B, HHH \\
XIV & DOMAIN-CATALOG & Ch. 1, 5 & All DOMAIN, X \\
BI & DOMAIN-LIFESPAN & Ch. 24 & Ch. 23, 17, RRR, OOO, PPP \\
\addlinespace

\multicolumn{4}{@{}l}{\textit{\textbf{CONTEXT (6)}}} \\
AH & CONTEXT-TIME & Ch. 9 & Ch. 10.7 \\
AI & CONTEXT-SPACE & Ch. 9 & Ch. 10.7 \\
BG & CONTEXT-GENESIS & Ch. 14 & BH, DX \\
BH & CONTEXT-SPÖ_PARADOX & Ch. 14 & BG, DX \\
DX & CONTEXT-MEDIA & Ch. 9, 14 & BG, BH, BS, BT, BF \\
V & CONTEXT-MASTER & Ch. 9 & Ch. 10.7, 11.15 \\
\addlinespace

\multicolumn{4}{@{}l}{\textit{\textbf{METHOD (12)}}} \\
AN & METHOD-LLMMC & Ch. 4x & Ch. 4 \\
AL & METHOD-SRL & Ch. 11 & Ch. 9 \\
E & METHOD-OPS & Ch. 4 & Ch. 10 \\
R & METHOD-EVAL & Ch. 4 & Ch. 18 \\
AZ & METHOD-CONSTRUCT & Ch. 8 & Ch. 2, 4x, 6 \\
CCC & METHOD-DOCTYPE & Ch. 17 & Ch. 4x, 9 \\
EEE & METHOD-DESIGN & Ch. 4 & Ch. 10, 18 \\
FFF & METHOD-REGISTRY & Ch. 4 & Ch. 10, 18 \\
GGG & METHOD-CONFIG & Ch. 4 & Ch. 9, 10 \\
HHH & METHOD-TOOLKIT & Ch. 17 & Ch. 15, HHH \\
VIII & METHOD-COMP & Ch. 5, 8 & B, VII, All FORMAL \\
IX & METHOD-RESEARCH & Ch. 18 & VII, VIII, All FORMAL \\
III & METHOD-LLMMC-CAL & Ch. 17 & HHH, EEE, BBB \\
XX & METHOD-EVIDENCE & Ch. 4 & BBB, All LIT, All FORMAL \\
BJ & METHOD-DOMAINVAL & Ch. 24 & BI, RRR, C \\
BK & METHOD-CALC & Ch. 24, 17 & AU, AV, AW, RRR, OOO, BI, BJ \\
BM & METHOD-PAPERINT & All & BBB, LRN, /upgrade-paper \\
VC & METHOD-VALIDATE & Ch. 9, 10 & V, BBB, FRM, AN \\
\addlinespace

\multicolumn{4}{@{}l}{\textit{\textbf{PREDICT (8)}}} \\
S & PREDICT-MASTER & Ch. 18 & Ch. 4 \\
AO & PREDICT-01 & Ch. 11.1 & Ch. 9 \\
AP & PREDICT-02 & Ch. 11.2 & Ch. 10.8 \\
AQ & PREDICT-03 & Ch. 11.3 & Ch. 10 \\
AR & PREDICT-04 & Ch. 11.4 & Ch. 9 \\
AS & PREDICT-05 & Ch. 11.5 & Ch. 10.1 \\
AT & PREDICT-06 & Ch. 11.6 & Ch. 10.8 \\
HF & PREDICT-HABIT & Ch. 13 & Ch. 14, AW, AU \\
\addlinespace

\multicolumn{4}{@{}l}{\textit{\textbf{LIT (16)}}} \\
I & LIT-NOBEL & Ch. 1 & All \\
J & LIT-RECENT & Ch. 4 & All \\
K & LIT-FEHR & Ch. 4 & Ch. 10.1 \\
L & LIT-ACEMOGLU & Ch. 9 & Ch. 10.8 \\
M & LIT-SHLEIFER & Ch. 3 & Ch. 5 \\
N & LIT-HECKMAN & Ch. 4 & Ch. 10.9 \\
O & LIT-AUTOR & Ch. 10.9.1 & Ch. 4 \\
P & LIT-DUFLO & Ch. 4 & Ch. 17 \\
Q & LIT-BLOOM & Ch. 9 & Ch. 5 \\
R & LIT-THALER & Ch. 14 & Ch. 16, 17 \\
S & LIT-SUNSTEIN & Ch. 14 & Ch. 16, 17 \\
T & LIT-CAMERER & Ch. 14 & Ch. 16, 17 \\
W & LIT-ARIELY & Ch. 14 & Ch. 16, 17 \\
X & LIT-LOEWENSTEIN & Ch. 14 & Ch. 16, 17 \\
Y & LIT-CIALDINI & Ch. 14 & Ch. 16, 17 \\
Z & LIT-HAIDT & Ch. 14 & Ch. 16, 17 \\
AA & LIT-MULLAINATHAN & Ch. 14 & Ch. 16, 17 \\
AB & LIT-LIST & Ch. 14 & Ch. 16, 17 \\
AC & LIT-DOLAN & Ch. 14 & Ch. 16, 17 \\
AD & LIT-SPECIALISTS & Ch. 14 & Ch. 16, 17 \\
AE & LIT-FALK & Ch. 5, 9 & Ch. 10, 15 \\
AF & LIT-MALMENDIER & Ch. 14 & Ch. 16, 17 \\
AG & LIT-SHAFIR & Ch. 14 & Ch. 16, 17 \\
SU & LIT-SUTTER & Ch. 9, 14 & Ch. 5, 10 \\
CA & LIT-CAMERER & Ch. 5, 9 & Ch. 12, 14 \\
LO & LIT-LOEWENSTEIN & Ch. 11, 9 & Ch. 13, AU \\
GI & LIT-GIGERENZER & Ch. 11, 3 & Ch. 9, L \\
RO & LIT-ROTH & Ch. 14, 5 & Ch. 9, AA, AJ \\
RA & LIT-RABIN & Ch. 6, 10 & Ch. 11, 12, U \\
FI & LIT-FISCHBACHER & Ch. 5, 9 & Ch. 10.1, K, AM \\
MM & LIT-MARECHAL & Ch. 10, 14 & Ch. 5, V, AH, K, FI \\
LA & LIT-LAIBSON & Ch. 11, 6 & Ch. 14, V, RA, R, AD \\
GN & LIT-GNEEZY & Ch. 9, 5 & Ch. 14, B, K, AB, AI \\
DV & LIT-DELLAVIGNA & Ch. 11, 14 & Ch. 9, LA, AB, RA, V \\
SA & LIT-AMBÜHL & Ch. 12, 15 & Ch. 9, AU, RO, GN, K \\
MN & LIT-NIEDERLE & Ch. 9, 14 & Ch. 5, GN, RO, SA, AI \\
XG & LIT-GABAIX & Ch. 11, 9 & Ch. 14, AU, LA, V, R \\
BE & LIT-ENKE & Ch. 9, 11 & Ch. 14, AU, V, AE, GN \\
PA & LIT-AGHION & Ch. 9, 14 & Ch. 5, L, Q, AC \\
DC & LIT-CARD & Ch. 4, 14 & Ch. 9, N, O, DV \\
GA & LIT-AKERLOF & Ch. 10, 14 & Ch. 5, K, V, AJ \\
VS & LIT-SMITH & Ch. 4, 9 & Ch. 5, 11, K, RO, CA \\
TS & LIT-SCHELLING & Ch. 5, 12 & Ch. 9, 14, CA, AV, AG \\
GB & LIT-BECKER & Ch. 3, 5, 7, 9 & Ch. 14, AA, AB, BBB \\
JT & LIT-TIROLE & Ch. 9, 10, 14 & IOO, AF, K, GB \\
RB & LIT-BENABOU & Ch. 5, 7, 10, 11 & JT, AU, AV, AB, AG \\
SE & LIT-STRATIFICATION & Ch. 7, 9, 14 & GB, RB, K, AG, AE \\
OT & LIT-OTHER & Variable & All LIT appendices \\
EV & LIT-VANDENSTEEN & Ch. 9, 10 & Ch. 5, 14, JT, RB, GA \\
RS & LIT-SADUN & Ch. 5, 9 & Ch. 17, Q, EV, PA \\
CG & LIT-GOLDIN & Ch. 9, 10 & Ch. 7, 14, GB, DC, N \\
RN & LIT-NAGEL & Ch. 11, 5 & Ch. 8, 12, CC, GG, K \\
U & LIT-KT & Ch. 3, 6 & Ch. 10 \\
AX & LIT-META & Ch. 1 & Ch. 18, AW \\
AY & LIT-PARADIGMS & Ch. 2 & Ch. 1, 6, 18 \\
XV & LIT-HISTORY & Ch. 1, 5 & VII, X, XIV \\
XVI & LIT-CRITIQUE & Ch. 5, 18 & XV, X, VII \\
XVII & LIT-THERMODYNAMICS & Ch. 1, 5 & B, XV, VII \\
XVIII & LIT-FEHR-METHOD & Ch. 5, 9 & \textbf{All CORE (10C)}, K, XVI, S \\
XIX & LIT-LEARNING & Ch. 5, 14 & HF, PBB, AU, BBB \\
\addlinespace

\multicolumn{4}{@{}l}{\textit{\textbf{REF (9)}}} \\
F & REF-EXAMPLES & Ch. 10.9 & All \\
G & REF-GLOSSARY & All & --- \\
H & REF-HISTORY & Ch. 1 & --- \\
T & REF-META & Ch. 2 & Ch. 18 \\
SK & REF-SKILLS & CLAUDE.md & All \\
DDD & REF-DOCTYPE & Ch. 17 & Ch. 9, 11 \\
DA & REF-DEFENSE & Ch. 1 & Ch. 18, All \\
DA2 & REF-DEFENSE-MATRIX & Ch. 1 & DA, All \\
BO & REF-ARCH & All & SK, G, BBB, All CORE \\
BN & REF-LEAD & N/A & G, LEAD-ID-SYSTEM.md \\

\end{longtable}

\subsection{Reading Pathways by Chapter}

\begin{tcolorbox}[colback=green!5!white, colframe=green!75!black, 
    title=How to Use: Chapter $\rightarrow$ Appendix]

\textbf{If you're reading Chapter 5 (Complementarity):}
\begin{itemize}[nosep]
\item Core theory: CORE-HOW (B)
\item Domain examples: DOMAIN-IO (AC), DOMAIN-MATCH (AB)
\item Foundational literature: DOMAIN-COMPLEMENT (X)
\end{itemize}

\textbf{If you're reading Chapter 9 (Context):}
\begin{itemize}[nosep]
\item Core theory: CORE-WHEN (V), CORE-WHO (AAA)
\item Temporal details: CONTEXT-TIME (AH)
\item Spatial details: CONTEXT-SPACE (AI)
\end{itemize}

\textbf{If you're reading Chapter 10 (Nutzenarchitektur):}
\begin{itemize}[nosep]
\item Complete specification: CORE-WHAT (C)
\item All axioms: FORMAL-AXIOM (AU)
\item Level structure: CORE-WHO (AAA)
\end{itemize}

\textbf{If you're reading Chapter 11 (Awareness):}
\begin{itemize}[nosep]
\item Willingness to act: CORE-READY (AV)
\item Predictions: PREDICT-01 through PREDICT-06
\item Level-specific awareness: CORE-WHO (AAA)
\end{itemize}

\textbf{If you're reading Chapter 17 (Intervention Foundations):}
\begin{itemize}[nosep]
\item Intervention toolkit: METHOD-TOOLKIT (HHH)
\item Efficient intervention: FORMAL-EIH (III)
\item Parameter sourcing: METHOD-PSG (JJJ)
\end{itemize}

\textbf{If you're reading Chapter 24 (Emergent Life Journeys):}
\begin{itemize}[nosep]
\item Intergenerational spillover: DOMAIN-LIFESPAN (BI)
\item Cross-cultural validation: METHOD-DOMAINVAL (BJ)
\item Calculation reference: METHOD-CALC (BK)
\item Lifetime parameters: METHOD-LIFETIME (RRR)
\item System dynamics: FORMAL-SYSTEM (OOO)
\item Agent-based modeling: METHOD-ABM (PPP)
\end{itemize}

\end{tcolorbox}

%%%%%%%%%%%%%%%%%%%%%%%%%%%%%%%%%%%%%%%%%%%%%%%%%%%%%%%%%%%%%%%%%%%%%%%%%%%%%%%
\section{Summary}
\label{sec:index-summary}
%%%%%%%%%%%%%%%%%%%%%%%%%%%%%%%%%%%%%%%%%%%%%%%%%%%%%%%%%%%%%%%%%%%%%%%%%%%%%%%

\begin{tcolorbox}[colback=gray!5!white, colframe=gray!75!black,
    title=Summary: The EBF Appendix System]

This index catalogs the complete EBF appendix ecosystem:

\begin{enumerate}[nosep]
\item \textbf{155 Appendices} organized into 8 categories (CORE, FORMAL, DOMAIN, CONTEXT, METHOD, PREDICT, LIT, REF)
\item \textbf{24 Main Chapters} with bidirectional linkages to supporting appendices
\item \textbf{10C Framework} at the core: WHO, WHAT, HOW, WHEN, WHERE, AWARE, READY, STAGE, HIERARCHY
\item \textbf{Reading Pathways} for different user types (theorist, empiricist, practitioner)
\end{enumerate}

Every appendix follows the naming convention \textbf{CATEGORY-NAME: Descriptive Title} and specifies its primary and secondary chapter linkages.

\end{tcolorbox}

%%%%%%%%%%%%%%%%%%%%%%%%%%%%%%%%%%%%%%%%%%%%%%%%%%%%%%%%%%%%%%%%%%%%%%%%%%%%%%%
\section{Glossary}
\label{sec:index-glossary}
%%%%%%%%%%%%%%%%%%%%%%%%%%%%%%%%%%%%%%%%%%%%%%%%%%%%%%%%%%%%%%%%%%%%%%%%%%%%%%%

For complete terminology definitions, see \textbf{Appendix G: REF-GLOSSARY} (\ref{app:G}).

Key terms used in this index:
\begin{itemize}[nosep]
\item \textbf{10C}: The ten CORE questions (WHO, WHAT, HOW, WHEN, WHERE, AWARE, READY, STAGE, HIERARCHY, plus EIT), with PRINCIPLES (PR) defining EBF membership criteria
\item \textbf{Primary Chapter}: The main conceptual home of an appendix
\item \textbf{Secondary Chapter}: Additional chapters that reference the appendix
\item \textbf{Reading Pathway}: Recommended sequence of appendices for specific goals
\end{itemize}

%%%%%%%%%%%%%%%%%%%%%%%%%%%%%%%%%%%%%%%%%%%%%%%%%%%%%%%%%%%%%%%%%%%%%%%%%%%%%%%
\section{References}
\label{sec:index-references}
%%%%%%%%%%%%%%%%%%%%%%%%%%%%%%%%%%%%%%%%%%%%%%%%%%%%%%%%%%%%%%%%%%%%%%%%%%%%%%%

This appendix index draws on the complete EBF bibliography. For the full reference list, see the Master Bibliography (bcm2\_master\_references.bib).

\nocite{bcm_master}

\begin{itemize}[nosep]
\item \textbf{Master Bibliography} (bcm2\_master\_references.bib): 1,922+ papers from behavioral economics research
\item \textbf{LIT Appendices}: Author-specific research integrations (63 appendices)
\item \textbf{Nobel Contributions (I)}: Historical foundations
\end{itemize}

\end{document}
